\documentclass[11pt]{article}
\usepackage[utf8]{inputenc}
\usepackage[margin=1in]{geometry}
\usepackage{enumitem}
\usepackage{titlesec}
\usepackage{xcolor}

\titleformat{\section}{\large\bfseries\color{blue!70!black}}{}{0em}{}[\titlerule]
\titleformat{\subsection}{\normalsize\bfseries}{}{0em}{}

\title{\textbf{Tabletop Exercise: The Recursive Trust Hijack}\\ \large Red-Teaming Identity-to-Control-Plane Escalation}
\author{Based on the Research of Rogel S.J. Corral}
\date{March 2026}

\begin{document}

\maketitle

\section{Exercise Overview}
This exercise validates an organization’s resilience against \textbf{Recursive Trust} chains[cite: 28, 96]. It tests the transition from a simple "trust error" into an "identity event," eventually targeting the administrative control plane[cite: 12, 34].

\section{The Scenario}

\subsection{Phase 1: Adaptive Reconnaissance \& The Lure}
An attacker uses an LLM to synthesize public data into a high-fidelity message matching a mid-level employee’s workflow[cite: 35, 37]. 
\begin{itemize}
    \item \textbf{Inject:} The target receives a DM or email regarding a routine "Budget Outlook" or "Access Review." [cite: 33]
    \item \textbf{The Breach:} The employee enters credentials into a spoofed "Secure Portal" flow. [cite: 39]
\end{itemize}

\subsection{Phase 2: Post-Compromise Expansion}
The attacker uses the initial session token to map privilege touchpoints, looking for delegation paths and group sprawl[cite: 43]. 
\begin{itemize}
    \item \textbf{The Pivot:} The attacker discovers "Delegation Drift"—a legacy permission grant that allows the user to influence sensitive policy objects[cite: 47, 94].
\end{itemize}

\subsection{Phase 3: Control-Plane Escalation}
The attacker converts shadow authority into administrative capability[cite: 48].
\begin{itemize}
    \item \textbf{The Impact:} A rapid cluster of changes occurs: a new "Break-Glass" account is created, or MFA is disabled for a specific administrative role[cite: 70].
\end{itemize}

\newpage

\section{One-Page Observer’s Guide}

\subsection{Objective}
To identify gaps in \textbf{Detection} (Time-to-Identify) and \textbf{Discipline} (Verification Rituals)[cite: 69, 89].

\subsection{Key Observation Points}
\begin{enumerate}
    \item \textbf{Authentication Ambiguity:} Did the participant question the "familiar" login prompt, or did familiarity substitute for verification? [cite: 41, 59]
    \item \textbf{Verification Rituals:} Did anyone suggest an out-of-band (OOB) check (e.g., a phone call or secondary chat) before complying with the request? [cite: 73, 81]
    \item \textbf{Detection Capability:} Would our SOC see a rapid cluster of policy edits following an identity event? [cite: 70, 84]
    \item \textbf{Blast Radius:} Is the compromised user’s access strictly limited by "Least Privilege," or does "Delegation Drift" allow lateral movement? [cite: 46, 66]
\end{enumerate}

\subsection{Success Metrics Tracking}
\begin{itemize}
    \item \textbf{Mean Time to Detect (MTTD):} Time elapsed between the "Identity Event" and SOC notification. [cite: 71]
    \item \textbf{Mean Time to Contain (MTTC):} Time elapsed before the administrative elevation was blocked. [cite: 71]
    \item \textbf{Verification Rate:} \% of participants who successfully triggered a verification ritual. [cite: 75]
\end{itemize}

\subsection{Debrief Questions}
\begin{itemize}
    \item Do we have a standard script for staff who feel pressured by "senior leadership" messages? [cite: 74]
    \item Are our reset and recovery workflows hardened with two-person integrity? [cite: 64]
    \item Is there a fast containment playbook for suspected identity-to-control-plane escalation? [cite: 85]
\end{itemize}

\end{document}
